% mathnote-colors.tex
% 屏幕 / 印刷双模式的颜色与调色板定义，
% 以及与 hyperref / PDF 元信息配色的联动。

% --------------------------------------------------
% 基础色
% --------------------------------------------------
\definecolor{screenAccent}{HTML}{1565C0}
\definecolor{screenSecondary}{HTML}{00897B}
\definecolor{screenHighlight}{HTML}{F9A826}
\definecolor{screenInfo}{HTML}{546E7A}
\definecolor{screenSurface}{HTML}{FFFFFF}

\definecolor{printAccent}{cmyk}{0.95,0.55,0,0.05}
\definecolor{printSecondary}{cmyk}{0.82,0,0.56,0.08}
\definecolor{printHighlight}{cmyk}{0,0.35,0.80,0}
\definecolor{printInfo}{cmyk}{0.60,0.47,0.43,0.20}
\definecolor{printSurface}{cmyk}{0,0,0,0}
\definecolor{mathnotePureCyan}{cmyk}{1,0,0,0}

% 16 色屏幕调色板（sRGB）
\definecolor{screenTone01}{HTML}{0D47A1}
\definecolor{screenTone02}{HTML}{1565C0}
\definecolor{screenTone03}{HTML}{1A73E8}
\definecolor{screenTone04}{HTML}{2196F3}
\definecolor{screenTone05}{HTML}{00ACC1}
\definecolor{screenTone06}{HTML}{00897B}
\definecolor{screenTone07}{HTML}{2E7D32}
\definecolor{screenTone08}{HTML}{558B2F}
\definecolor{screenTone09}{HTML}{9E9D24}
\definecolor{screenTone10}{HTML}{F9A825}
\definecolor{screenTone11}{HTML}{FFB300}
\definecolor{screenTone12}{HTML}{FB8C00}
\definecolor{screenTone13}{HTML}{F4511E}
\definecolor{screenTone14}{HTML}{D84315}
\definecolor{screenTone15}{HTML}{8E24AA}
\definecolor{screenTone16}{HTML}{6A1B9A}

% 16 色印刷调色板（CMYK，适合激光/数码印刷）
\definecolor{printTone01}{cmyk}{1,0.72,0,0.35}
\definecolor{printTone02}{cmyk}{0.9,0.5,0,0.2}
\definecolor{printTone03}{cmyk}{0.85,0.45,0,0.12}
\definecolor{printTone04}{cmyk}{0.65,0.25,0,0.02}
\definecolor{printTone05}{cmyk}{0.75,0.05,0.1,0.05}
\definecolor{printTone06}{cmyk}{0.85,0,0.35,0.2}
\definecolor{printTone07}{cmyk}{0.75,0,0.8,0.38}
\definecolor{printTone08}{cmyk}{0.6,0,1,0.42}
\definecolor{printTone09}{cmyk}{0.35,0,1,0.45}
\definecolor{printTone10}{cmyk}{0,0.2,1,0.02}
\definecolor{printTone11}{cmyk}{0,0.25,1,0}
\definecolor{printTone12}{cmyk}{0,0.45,1,0}
\definecolor{printTone13}{cmyk}{0,0.7,0.8,0}
\definecolor{printTone14}{cmyk}{0,0.85,0.95,0.1}
\definecolor{printTone15}{cmyk}{0.45,0.9,0,0}
\definecolor{printTone16}{cmyk}{0.6,1,0,0.1}

\newif\ifmathnote@docstarted
\mathnote@docstartedfalse
\AtBeginDocument{\mathnote@docstartedtrue}

% --------------------------------------------------
% hyperref 颜色联动
% --------------------------------------------------
\newcommand{\mathnote@applyhypercolors}{%
  \ifmathnoteprintmode
    \hypersetup{%
      colorlinks=false,
      hidelinks,
      pdfborderstyle={/S/U/W 0},
      pdfborder={0 0 0}%
    }%
  \else
    \hypersetup{%
      colorlinks=true,
      linkcolor=accent,
      citecolor=secondary,
      urlcolor=accent,
      pdfborder={0 0 0}%
    }%
  \fi
}

% --------------------------------------------------
% 主调色逻辑（基础语义色）
% --------------------------------------------------
\newcommand{\mathnote@applypalette}{%
  \ifmathnoteprintmode
    % 基础语义色（印刷 CMYK）
    \colorlet{accent}{printAccent}%
    \colorlet{secondary}{printSecondary}%
    \colorlet{highlight}{printHighlight}%
    \colorlet{inkgray}{printInfo}%
    \colorlet{surface}{printSurface}%
  \else
    % 基础语义色（屏幕 sRGB）
    \colorlet{accent}{screenAccent}%
    \colorlet{secondary}{screenSecondary}%
    \colorlet{highlight}{screenHighlight}%
    \colorlet{inkgray}{screenInfo}%
    \colorlet{surface}{screenSurface}%
  \fi
  % 盒子和网格的派生色
  \colorlet{accentline}{accent!65!black}%
  \colorlet{accentbg}{accent!8!white}%
  \colorlet{secondarybg}{secondary!10!white}%
  \colorlet{highlightbg}{highlight!10!white}%
  \colorlet{inkline}{inkgray!60!black}%
  % colframe Ϊ tcolorbox �Ļ�����ɫ��Ϊ�������� coverage=soft/solid
  % �����漰 colframe!x!white ��ɫ����ʧ�ҳ���һ��ȱʡֵ
  \colorlet{colframe}{accentline}%
  \colorlet{surfacegrid}{inkgray!6!white}%
  \ifmathnote@docstarted
    \mathnote@applyhypercolors
  \else
    \AtBeginDocument{\mathnote@applyhypercolors}
  \fi
}

% --------------------------------------------------
% 高层接口：16 色调色板 + 标题用色
% --------------------------------------------------
\newcommand{\MathNoteApplyPalette}{%
  \mathnote@applypalette
  % 16 色预置调色板：tone01–tone16 统一语义名，
  % 根据屏幕/印刷模式自动映射到 RGB / CMYK。
  \ifmathnoteprintmode
    \colorlet{tone01}{printTone01}%
    \colorlet{tone02}{printTone02}%
    \colorlet{tone03}{printTone03}%
    \colorlet{tone04}{printTone04}%
    \colorlet{tone05}{printTone05}%
    \colorlet{tone06}{printTone06}%
    \colorlet{tone07}{printTone07}%
    \colorlet{tone08}{printTone08}%
    \colorlet{tone09}{printTone09}%
    \colorlet{tone10}{printTone10}%
    \colorlet{tone11}{printTone11}%
    \colorlet{tone12}{printTone12}%
    \colorlet{tone13}{printTone13}%
    \colorlet{tone14}{printTone14}%
    \colorlet{tone15}{printTone15}%
    \colorlet{tone16}{printTone16}%
  \else
    \colorlet{tone01}{screenTone01}%
    \colorlet{tone02}{screenTone02}%
    \colorlet{tone03}{screenTone03}%
    \colorlet{tone04}{screenTone04}%
    \colorlet{tone05}{screenTone05}%
    \colorlet{tone06}{screenTone06}%
    \colorlet{tone07}{screenTone07}%
    \colorlet{tone08}{screenTone08}%
    \colorlet{tone09}{screenTone09}%
    \colorlet{tone10}{screenTone10}%
    \colorlet{tone11}{screenTone11}%
    \colorlet{tone12}{screenTone12}%
    \colorlet{tone13}{screenTone13}%
    \colorlet{tone14}{screenTone14}%
    \colorlet{tone15}{screenTone15}%
    \colorlet{tone16}{screenTone16}%
  \fi
  % 标题文字颜色：略降饱和度，更“克制”
  \colorlet{accenttitle}{accent!25!black}%
  \colorlet{secondarytitle}{secondary!25!black}%
  \colorlet{highlighttitle}{highlight!25!black}%
  \colorlet{inktitle}{inkgray!25!black}%
}

% 初始化为屏幕/印刷对应的完整调色板
\MathNoteApplyPalette
% 用户在外部修改 \colorlet 后，可显式调用以刷新
\newcommand{\MathNoteRefreshColors}{\MathNoteApplyPalette}

% PDF 元信息颜色
\AtBeginDocument{%
  \hypersetup{%
    pdftitle=\notetitle,
    pdfauthor=\noteauthor,
    pdfsubject=\notesubtitle,
    pdfcreator={MathNote dual-medium template}%
  }%
}
